% Applied Software Development Lab -- Lab Handout 3: Blinking LED -- Bare-Metal C
% Copyright (c) Michael Lehn. Licensed under CC BY-SA 4.0.
%
% Region coordinates (bp from the top of the page) follow the page layout of
% the original German handout "Laborblatt 3".
\documentclass{aslab}
\sethandout{3}{Blinking LED -- Bare-Metal C}
\graphicspath{{../figures/handout03/}{../figures/shared/}}

\makeatletter
%% Lists used in this handout --------------------------------------------
%% learning objectives: blue bold dashes, text at +12bp
\newlist{bluedash}{itemize}{1}
\setlist[bluedash]{label={\color{aslabblue}\bfseries\aslab@dash},
  leftmargin=12bp,labelsep=*,labelindent=0pt,align=left}
%% "A)" lists with the text at +18bp
\newlist{alphatasks}{enumerate}{1}
\setlist[alphatasks]{label={\Alph*)},leftmargin=18bp,labelsep=*,align=left,
  labelindent=0pt}
%% "a)" lists with the text at +18bp
\newlist{alphasub}{enumerate}{1}
\setlist[alphasub]{label={\alph*)},leftmargin=18bp,labelsep=*,align=left,
  labelindent=0pt}
%% plain "1." lists (regular weight)
\newlist{numlist}{enumerate}{1}
\setlist[numlist]{label={\arabic*.},leftmargin=18bp,labelsep=*,align=left,
  labelindent=0pt}

%% Assembly and C listings with syntax colors ----------------------------
%% Inside asmlisting, every space is a fixed 0.6em and every line break
%% ends a line, so the columns of the original listings are preserved.
%% The optional argument is the font size in bp (default 11); the line
%% pitch is the one of the surrounding region.
\newcommand*\aN[1]{\textcolor{codeorange}{#1}}% line numbers, labels, keywords
\newcommand*\aI[1]{\textcolor{codegreen}{#1}}% instructions, types
\newcommand*\aK[1]{\textcolor{codered}{#1}}% constants, I/O addresses
\newcommand*\aR[1]{\textcolor{codecyan}{#1}}% registers, jump targets
\newcommand*\aC[1]{\textcolor{codepurple}{#1}}% comments
\newcommand*\asm@space{\leavevmode\hbox to 0.6em{}}
{\catcode`\ =\active\catcode`\^^M=\active%
\gdef\asm@setup{\catcode`\ =\active\def {\asm@space}%
\catcode`\^^M=\active\def^^M{\par}}}%
\newenvironment{asmlisting}[1][11]{\par
  \edef\asm@bs{\the\baselineskip}\fontsize{#1bp}{\asm@bs}\selectfont
  \ttfamily\asm@setup}{\par}

%% Keep the baseline pitch fixed even next to tall inline code boxes
\lineskiplimit=-\maxdimen
\makeatother

\begin{document}
%% ===================================================================== 1 ==
\maketitleblock

\begin{region}{124.0}
\heading{Learning objectives}
\end{region}

\begin{region}{148.5}
\begin{bluedash}
\item Introduction to programming in C/C++ instead of assembly language
\item Interpreting and analyzing the assembly code generated by the C
  compiler
\item Using the oscilloscope to analyze digital output signals
\end{bluedash}
\end{region}

\begin{region}{199.0}
\heading{Step 1: Working with the terminal and experimental setup}
\end{region}

\begin{region}[pitch=13.25]{227.5}
\textbf{Goal:} Building on the last lab session: we recap how we edit
programs, compile them, and transfer them to the microcontroller. We also
check the blinking frequency of the LED with the oscilloscope.
\end{region}

\begin{region}[pitch=13.1]{280.4}
As usual, we work with two terminals side by side -- the left one for
editing, the right one for running commands.
\end{region}

\begin{region}{319.6}
\subheading{Left terminal (editing)}
\end{region}

\begin{region}[pitch=12.17]{333.9}
\begin{alphatasks}
\item \textbf{Clone the Git repository}\\
  Download the prepared code base with the following command:\\
  \code{git clone https://gitlab.uni-ulm.de/mlehn/avr-bare-metal-2.git}\\
  Among other things, the repository contains the file
  \code{led_pin2_delay4.s}, an assembly program that makes an LED connected
  to pin~2 blink at 1\,Hz.
\end{alphatasks}
\end{region}

\begin{region}[pitch=12.15]{395.5}
\begin{alphatasks}[start=2]
\item \textbf{Change into the subdirectory and open the file in nano}\\
  Change into the directory with the assembly code:\\
  \code{cd avr-bare-metal-2/asm} and open \code{led_pin2_delay4.s} in
  \code{nano}.
\end{alphatasks}
\end{region}

\begin{region}{444.8}
\subheading{Right terminal (compiling \& flashing)}
\end{region}

\begin{region}[pitch=12.2]{457.8}
\begin{alphatasks}
\item \textbf{Change into the same directory and flash}\\
  Here, too, first change into the subdirectory:\\
  \code{cd avr-bare-metal-2/asm}\\
  Connect your microcontroller and make it visible to WSL with
  \code{usb_check}. Then assemble and upload the code with:\\
  \code{make upload HEX=led_pin2_delay4.hex}
\end{alphatasks}
\end{region}

\begin{region}{543.4}
\subheading{Task}
\end{region}

\begin{region}{567.5}
\begin{numlist}
\item Connect an LED with a suitable series resistor to pin~2 and watch it
  blink.
\item Measure with the oscilloscope at which frequency the LED actually
  blinks.
\item Explain and demonstrate at least the following points:
  \begin{steps}
  \item How does the assembly program make pin~2 oscillate at 1\,Hz?
  \item Which tool translates the assembly code into machine code?
  \item What does the machine code that is flashed look like?
  \item Which tool is used to flash the microcontroller?
  \end{steps}
\end{numlist}
\end{region}

\nextpage
%% ===================================================================== 2 ==
\begin{region}{69.4}
\heading{Step 2: Modifications to the assembly program}
\end{region}

\begin{region}[pitch=13.17]{83.5}
Investigate how different assembly instructions affect the behavior of the
program -- in particular the frequency at pin~2.\\
You will modify the program step by step and check each time with the
oscilloscope whether the LED still blinks at 1\,Hz.
\end{region}

\begin{region}{139.3}
\subheading{Why this matters:}
\end{region}

\begin{region}{152.6}
\begin{dashlist}
\item \textbf{Officially:} You will get to know more instructions of the AVR
  architecture.
\end{dashlist}
\end{region}

\begin{region}[pitch=14.3]{166.9}
\begin{dashlist}
\item \textbf{Unofficially:} You will create variants of assembly code such
  as a C compiler could also produce. This way, you will later see with
  your own eyes:
\end{dashlist}
\end{region}

\begin{region}[x=68.7,right=0]{194.3}
\begin{dashlist}
\item \textbf{C code can be translated into different assembly code.}
\end{dashlist}
\end{region}

\begin{region}[x=68.7,right=0]{208.6}
\begin{dashlist}
\item \textbf{Assembly code is translated into machine code} -- and
  unambiguously so.
\end{dashlist}
\end{region}

\begin{region}[x=68.7,right=0]{222.4}
\begin{dashlist}
\item Conclusion: C is more convenient, but less precise than assembly
  language.
\end{dashlist}
\end{region}

\begin{region}[pitch=10]{239.9}
\begin{asmlisting}[9]
  \aN{1}         \aI{ldi}     \aR{r19},    \aK{4}           \aC{; for pin 2 (4 = 2\textasciicircum{}2)}
  \aN{2}         \aI{out}     \aK{0x0a},   \aR{r19}
  \aN{3}         \aI{ldi}     \aR{r18},    \aK{0}
  \aN{4 loop}:   \aI{eor}     \aR{r18}\aN{,}    \aR{r19}         \aC{; 1 cycle}
  \aN{5}         \aI{out}     \aK{0x0b},   \aR{r18}         \aC{; 1 cycle}
  \aN{6}         \aI{ldi}     \aR{r24},    \aK{0x54}        \aC{; 1 cycle}
  \aN{7}         \aI{ldi}     \aR{r25},    \aK{0x58}        \aC{; 1 cycle}
  \aN{8}         \aI{ldi}     \aR{r26},    \aK{0x14}        \aC{; 1 cycle}
  \aN{9}         \aI{ldi}     \aR{r27},    \aK{0x00}        \aC{; 1 cycle}
 \aN{10 delay}:  \aI{subi}    \aR{r24},    \aK{1}           \aC{; 1 cycle}
 \aN{11}         \aI{sbci}    \aR{r25},    \aK{0}           \aC{; 1 cycle}
 \aN{12}         \aI{sbci}    \aR{r26},    \aK{0}           \aC{; 1 cycle}
 \aN{13}         \aI{sbci}    \aR{r27},    \aK{0}           \aC{; 1 cycle}
 \aN{14}         \aI{brne}    \aR{delay}               \aC{; 2 cycles if taken, 1 otherwise}
 \aN{15}         \aI{rjmp}    \aR{loop}                \aC{; 2 cycles}
\end{asmlisting}
\end{region}

\begin{region}{398.2}
\subheading{Tasks (measure with the oscilloscope after each step)}
\end{region}

\begin{region}{414.2}
\begin{numlist}
\item Replace the following two lines:
\end{numlist}
\end{region}

\begin{region}[x=74.7,right=0,pitch=10]{424.9}
\begin{asmlisting}[9]
\aN{10 delay}\aC{:}  \aI{subi}    \aR{r24},    \aK{1}           \aC{; 1 cycle}
\aN{11}         \aI{sbci}    \aR{r25},    \aK{0}           \aC{; 1 cycle}
\end{asmlisting}
\end{region}

\begin{region}[x=74.7,right=0]{446.2}
with this single line:
\end{region}

\begin{region}[x=74.7,right=0]{456.9}
\begin{asmlisting}[9]
\aN{10}         \aI{sbiw}    \aR{r24},    \aK{1}           \aC{; 2 cycles}
\end{asmlisting}
\end{region}

\begin{region}[x=74.7,right=0,pitch=12.03]{469.2}
The instruction \code{sbiw} \textit{(Subtract Immediate from Word)} treats
the register pair r25:r24 as a 16-bit value and subtracts 1 here. Although
it is not faster (also 2 cycles), the code becomes 2 bytes smaller --
because sbiw occupies only one instruction (2 bytes) instead of two
(2\,×\,2\,bytes).
\end{region}

\begin{region}{519.3}
\begin{numlist}[start=2]
\item Replace the following two lines:
\end{numlist}
\end{region}

\begin{region}[x=74.7,right=0,pitch=10]{530.0}
\begin{asmlisting}[9]
\aN{12}         \aI{sbci}    \aR{r26},    \aK{0}           \aC{; 1 cycle}
\aN{13}         \aI{sbci}    \aR{r27},    \aK{0}           \aC{; 1 cycle}
\end{asmlisting}
\end{region}

\begin{region}[x=74.7,right=0]{551.3}
with:
\end{region}

\begin{region}[x=74.7,right=0,pitch=10.8]{562.0}
\begin{asmlisting}[9]
\aN{12}         \aI{sbc}     \aR{r26},    \aR{r1}          \aC{; 1 cycle}
\aN{13}         \aI{sbc}     \aR{r27},    \aR{r1}          \aC{; 1 cycle}
\end{asmlisting}
\end{region}

\begin{region}[x=74.7,right=0,pitch=12.47]{585.3}
The instruction \code{sbc} \textit{(Subtract with Carry)} works like
\code{sbci} \textit{(Subtract Immediate with Carry)}, except that the second
operand here is a register instead of an immediate value. Register \code{r1}
is automatically initialized with 0 on every reset. As long as you do not
change it yourself, it therefore always contains the value 0.
\end{region}

\begin{region}[x=74.7,right=0,pitch=12.1]{635.6}
\textbf{Conclusion:} Functionally, and in terms of length and clock cycles,
the code is identical, but you learn an alternative way of subtracting the
carry bit.
\end{region}

\begin{region}{661.7}
\begin{numlist}[start=3]
\item Replace the following three lines:
\end{numlist}
\end{region}

\begin{region}[x=74.7,right=0,pitch=10]{672.4}
\begin{asmlisting}[9]
\aN{1}         \aI{ldi}     \aR{r19},    \aK{4}           \aC{; for pin 2 (4 = 2\textasciicircum{}2)}
\aN{2}         \aI{out}     \aK{0x0a},   \aR{r19}
\aN{3}         \aI{ldi}     \aR{r18},    \aK{0}
\end{asmlisting}
\end{region}

\begin{region}[x=74.7,right=0]{703.7}
with:
\end{region}

\begin{region}[x=74.7,right=0,pitch=10]{714.4}
\begin{asmlisting}[9]
\aN{1}         \aI{ldi}     \aR{r24},    \aK{4}           \aC{; for pin 2 (4 = 2\textasciicircum{}2)}
\aN{2}         \aI{out}     \aK{0x0a},   \aR{r24}
\aN{3}         \aI{ldi}     \aR{r18},    \aK{0}
\aN{4}         \aI{ldi}     \aR{r19},    \aK{4}
\end{asmlisting}
\end{region}

\begin{region}[x=74.7,right=0]{755.7}
The program works as before -- explain why this change is nevertheless
not really a good idea.
\end{region}

\nextpage
%% ===================================================================== 3 ==
\begin{region}[x=57.9,right=0]{69.4}
\heading{Step 3: Bare-metal C}
\end{region}

\begin{region}[pitch=13.4]{84.1}
In the subdirectory \code{c} of the repository, you will find C programs for
the microcontroller.\\
To get there from the current directory \code{asm}, you have to go ``one
directory up and then one down'':
\end{region}

\begin{region}[pitch=13.2]{124.3}
\code{cd ..}\\
\code{cd c}
\end{region}

\begin{region}{151.0}
Or directly in one go with: \code{cd ../c}
\end{region}

\begin{region}{166.7}
In the left terminal, open the file \code{led_pin2_simple.c}:
\end{region}

\begin{region}[pitch=10]{180.9}
\begin{asmlisting}[9]
  \aN{1} \aI{static} \aI{void} setPinModes(\aI{unsigned} \aI{char} mode)
  \aN{2} \{
  \aN{3}     *(\aI{volatile} \aI{unsigned} \aI{char} *)(\aK{0x2A}) = mode;
  \aN{4} \}
  \aN{5}
  \aN{6} \aI{static} \aI{void} digitalWriteAll(\aI{unsigned} \aI{char} val)
  \aN{7} \{
  \aN{8}     *(\aI{volatile} \aI{unsigned} \aI{char} *)(\aK{0x2B}) = val;
  \aN{9} \}
 \aN{10}
 \aN{11} \aI{static} \aI{void} delay\_cycles(\aI{unsigned} \aI{long} count)
 \aN{12} \{
 \aN{13}     \aN{do} \{
 \aN{14}         \aC{// the compiler should use a register for count}
 \aN{15}         \aN{asm} \aI{volatile}(\aK{""} : \aK{"+r"}(count));
 \aN{16}         count = count - \aK{1};
 \aN{17}     \} \aN{while} (count != \aK{0});
 \aN{18} \}
 \aN{19}
 \aN{20} \aI{int} main()
 \aN{21} \{
 \aN{22}     setPinModes(\aK{4});
 \aN{23}     digitalWriteAll(\aK{4});
 \aN{24}     while (\aK{1}) \{
 \aN{25}     \}
 \aN{26} \}
\end{asmlisting}
\end{region}

\begin{region}{445.3}
\subheading{A few words about what you see there:}
\end{region}

\begin{region}[pitch=13.37]{458.1}
\begin{dashlist}
\item In the upper part (lines 1 to 18), three functions are defined:\\
  \code{setPinModes}, \code{digitalWriteAll}, and \code{delay_cycles}.\\
  You do not need to understand the implementation of these functions in
  detail -- at first glance, it may even look daunting, even to
  hard-boiled C programmers. But of course I am happy to explain the
  details to you if you are curious.
\end{dashlist}
\end{region}

\begin{region}[pitch=13.36]{524.5}
\begin{dashlist}
\item In the lower part (lines 20 to 26), you see the function \code{main}.
  When the program starts, the statements inside the curly braces are
  executed one after the other from top to bottom.\\
  Here, for example, \code{setPinModes(4)} and \code{digitalWriteAll(4)} are
  called -- functions with which the pins 0 to 7 can be manipulated.\\
  What specifically happens:
\end{dashlist}
\end{region}

\begin{region}[x=68.7,right=0,pitch=14.1]{604.2}
\begin{dashlist}
\item \code{setPinModes(4)} sets the data direction register so that only
  pin~2 is configured as an output. (4 corresponds to \code{00000100} in
  binary, i.e., only bit~2 is set.)
\end{dashlist}
\end{region}

\begin{region}[x=68.7,right=0,pitch=13.8]{631.2}
\begin{dashlist}
\item \code{digitalWriteAll(4)} then sets pin~2 HIGH, because here, too,
  only bit~2 is set in the argument. The other pins remain LOW.
\end{dashlist}
\end{region}

\begin{region}[x=68.7,right=0,pitch=13.5]{657.9}
\begin{dashlist}
\item After the output pin has been set with \code{digitalWriteAll(4)}, the
  microcontroller should simply do nothing else -- it should keep the state
  permanently.\\
  The loop
\end{dashlist}
\end{region}

\begin{region}[x=80.7,right=0,pitch=11]{698.9}
\begin{asmlisting}[9]
\aN{24}     while (\aK{1}) \{
\aN{25}     \}
\end{asmlisting}
\end{region}

\begin{region}[x=80.7,right=0,pitch=13.3]{724.4}
makes sure that the program does not end but stays in an endless loop. This
is common for microcontrollers: a program usually runs as long as the device
is powered -- and \code{main()} should therefore never return, i.e., never
terminate.
\end{region}

\nextpage
%% ===================================================================== 4 ==
\begin{region}{69.4}
\subheading{Task}
\end{region}

\begin{region}[pitch=13.3]{98.0}
\begin{numlist}
\item \textbf{Compiling and flashing}\\
  This directory (c) also contains a Makefile that can be used as usual. In
  the right terminal, run the following command to compile the program and
  flash it to the microcontroller:\\ \code{make upload HEX=led_pin2_simple.hex}
\end{numlist}
\end{region}

\begin{region}[x=74.7,right=0,pitch=13.1]{151.2}
\begin{alphasub}
\item Connect an LED to pin~2 via a suitable series resistor and make sure
  that it lights up.
\end{alphasub}
\end{region}

\begin{region}[x=74.7,right=0]{177.7}
\begin{alphasub}[start=2]
\item Pay attention to the output of \code{make}. Explain to us what has
  changed in the toolchain.
\end{alphasub}
\end{region}

\begin{region}[pitch=13.4]{191.9}
\begin{numlist}[start=2]
\item \textbf{Analyzing automatically generated assembly code}\\
  Use \code{make led_pin2_simple.s} to translate the C program into assembly
  code. Display its content with \code{cat led_pin2_simple.s}. The generated
  file contains a lot of additional information (e.g., file name, compiler
  version, comments) that is not relevant for us. Only a few lines are
  relevant, for example these:
\end{numlist}
\end{region}

\begin{region}[x=74.7,right=0,pitch=11]{257.3}
\begin{asmlisting}[9]
         \aI{ldi}     \aR{r24},    \aR{lo}(\aK{4})
         \aI{out}     \aK{0xa},    \aR{r24}
         \aI{out}     \aK{0xb},    \aR{r24}
\aN{.L2}:
         \aI{rjmp}    \aN{.L2}
\end{asmlisting}
\end{region}

\begin{region}[x=74.7,right=0]{313.6}
(The lo(4) here simply stands for the constant 4. You could also write
ldi r24, 4 directly.)
\end{region}

\begin{region}[pitch=13.4]{327.7}
\begin{numlist}[start=3]
\item \textbf{Writing and comparing your own version}\\
  In the left terminal, create the file check.s with the following content:
\end{numlist}
\end{region}

\begin{region}[x=74.7,right=0,pitch=11]{352.9}
\begin{asmlisting}[9]
\aN{1}         \aI{ldi}     \aR{r24},    \aK{4}
\aN{2}         \aI{out}     \aK{0xa},    \aR{r24}
\aN{3}         \aI{out}     \aK{0xb},    \aR{r24}
\aN{4 halt:}   \aI{rjmp}    \aR{halt}
\end{asmlisting}
\end{region}

\begin{region}[x=74.7,right=0]{411.7}
Generate machine code from it with \code{make check.hex}
\end{region}

\begin{region}[x=74.7,right=0,pitch=13.43]{424.9}
\begin{alphasub}
\item Compare the generated machine code with that of the C program:\\
  \code{cat led_pin2_simple.hex}\\
  \code{cat check.hex}\\
  Or use the \code{diff} command: \code{diff check.hex led_pin2_simple.hex}\\
  If no differences are found, the output stays empty -- otherwise
  \code{diff} shows the lines that differ.
\end{alphasub}
\end{region}

\begin{region}[x=74.7,right=0]{505.5}
\begin{alphasub}[start=2]
\item In \code{check.s}, change the first line:
\end{alphasub}
\end{region}

\begin{region}[x=92.7,right=0]{518.8}
\begin{asmlisting}[9]
\aN{1}         \aI{ldi}     \aR{r24},    \aK{4}
\end{asmlisting}
\end{region}

\begin{region}[x=92.7,right=0]{532.0}
to
\end{region}

\begin{region}[x=92.7,right=0]{545.1}
\begin{asmlisting}[9]
\aN{1}         \aI{ldi}     \aR{r24},    \aK{5}
\end{asmlisting}
\end{region}

\begin{region}[x=92.7,right=0,pitch=13.25]{558.5}
Assemble it again with \code{make check.hex} and compare again with the
\code{diff} command. Now a difference should be displayed -- check where the
generated .hex files differ.
\end{region}

\nextpage
%% ===================================================================== 5 ==
\begin{region}{69.4}
\heading{Step 4: Blinking more elegantly with exclusive OR}
\end{region}

\begin{region}{85.2}
Modify the function \code{main()} as follows:
\end{region}

\begin{region}[pitch=10]{97.2}
\begin{asmlisting}[9]
 \aN{20} \aI{int} main()
 \aN{21} \{
 \aN{22}     setPinModes(\aK{4});
 \aN{23}     \aI{unsigned} \aI{char} pinValues = \aK{4};
 \aN{24}     \aN{while} (\aK{1}) \{
 \aN{25}         pinValues = pinValues \textasciicircum{} \aK{4};
 \aN{26}         digitalWriteAll(pinValues);
 \aN{27}     \}
 \aN{28} \}
\end{asmlisting}
\end{region}

\begin{region}{196.4}
\subheading{Tasks}
\end{region}

\begin{region}[pitch=13.27]{217.7}
\begin{numlist}
\item \textbf{Testing the modified program}\\
  Before we analyze in more detail why the program works, we first test
  whether it works:
\end{numlist}
\end{region}

\begin{region}[x=74.7,right=0,pitch=13.5]{264.3}
\begin{bullets}
\item Compile the program with the Makefile as usual and flash it to the
  microcontroller\\ with: \code{make upload HEX=led_pin2_simple.hex}
\end{bullets}
\end{region}

\begin{region}[x=74.7,right=0]{298.0}
\begin{bullets}
\item Observe what the LED at pin~2 does.
\end{bullets}
\end{region}

\begin{region}[x=74.7,right=0]{318.2}
\begin{bullets}
\item Measure with the oscilloscope how the signal at pin~2 behaves.
\end{bullets}
\end{region}

\begin{region}[pitch=13.43]{332.3}
\begin{numlist}[start=2]
\item \textbf{Writing and comparing your own version}\\
  Compile the C program with \code{make led_pin2_simple.s}. This generates
  the assembly code produced by the compiler. Display its content with:
  \code{cat led_pin2_simple.s}\\
  The part of the generated code that is relevant for us is equivalent to:
\end{numlist}
\end{region}

\begin{region}[x=74.7,right=0,pitch=11]{397.7}
\begin{asmlisting}[9]
\aN{1}         \aI{ldi}     \aR{r24},    \aK{4}
\aN{2}         \aI{out}     \aK{0xa},    \aR{r24}
\aN{3}         \aI{ldi}     \aR{r25},    \aK{4}
\aN{4 loop}:   \aI{eor}     \aR{r24},    \aR{r25}
\aN{5}         \aI{out}     \aK{0xb},    \aR{r24}
\aN{6}         \aI{rjmp}    \aR{loop}
\end{asmlisting}
\end{region}

\begin{region}[x=74.7,right=0,pitch=13.5]{478.4}
In the left terminal, create the file \code{check.s} with exactly this
content. Build both programs with: \code{make led_pin2_simple.hex check.hex}
and compare the generated machine code with \code{diff}.
\end{region}

\begin{region}{525.1}
\subheading{Brief explanation of the C code}
\end{region}

\begin{region}{542.8}
In line 23,
\end{region}

\begin{region}{557.0}
\begin{asmlisting}[9]
\aN{23}     \aI{unsigned} \aI{char} pinValues = \aK{4};
\end{asmlisting}
\end{region}

\begin{region}{572.3}
a \textbf{local} variable of type \code{unsigned char} is defined and
initialized with 4:
\end{region}

\begin{region}{588.4}
\begin{bullets}
\item The type \code{unsigned char} specifies that the variable occupies
  exactly one byte.
\end{bullets}
\end{region}

\begin{region}[pitch=13.84]{601.9}
\begin{bullets}
\item Since pinValues was defined inside the function \code{main} (without
  special additional keywords), it is a local variable with automatic
  storage duration: it exists only while the function is executing and is
  no longer valid after the function has ended.\\
  The compiler may place such \textbf{local} variables in registers, which
  is particularly efficient. If, on the other hand, a variable is defined
  outside all functions, it is \textbf{global} and permanently requires
  memory in RAM.
\end{bullets}
\end{region}

\begin{region}{685.4}
In line 25,
\end{region}

\begin{region}{701.2}
\begin{asmlisting}[9]
\aN{25}         pinValues = pinValues \textasciicircum{} \aK{4};
\end{asmlisting}
\end{region}

\begin{region}[pitch=15.5]{716.5}
the current value of pinValues is replaced by its exclusive OR with 4.\\
As you know, this has the effect that bit~2 is toggled (flipped):
\end{region}

\begin{region}[pitch=17.2]{747.7}
\begin{dashlist}
\item if bit~2 was 1 before $\rightarrow$ it becomes 0,
\item if it was 0 $\rightarrow$ it becomes 1.
\end{dashlist}
\end{region}

\begin{region}{782.0}
This creates the blinking effect.
\end{region}

\nextpage
%% ===================================================================== 6 ==
\begin{region}{69.4}
\heading{Step 5: Calmer blinking at 1\,Hz}
\end{region}

\begin{region}[pitch=14.3]{98.0}
Great -- once again we have a program that switches \textbf{pin~2
alternately HIGH for about 250\,ns and then LOW}.
\end{region}

\begin{region}[pitch=13.8]{126.5}
This time, however, the program is \textbf{written in C}, translated by the
\textbf{C compiler into assembly code}, and then by the \textbf{assembler
into machine code}. So, in a roundabout way, we have achieved the same as
before -- but with more modern means.
\end{region}

\begin{region}[pitch=13.75]{180.4}
Even if the C code (even main() alone) looks cryptic to a layperson at
first, you can sense that it is considerably \textbf{easier to learn the
rules of C} than to program directly in assembly language or even machine
code -- especially for larger projects.
\end{region}

\begin{region}[pitch=13.17]{234.2}
In the last example, the LED should blink at 1\,Hz again. To do this, we
insert a delay at the right place in the loop that waits for half a second
each time -- this results in a complete cycle of one second (1\,Hz). The
required number of clock cycles can be calculated:
\end{region}

\begin{region}[x=92.7,right=0]{303.6}
\fontsize{13.2bp}{14bp}\selectfont
$n = \dfrac{\frac{500000000}{62.5} - 7}{6} = 1333332$
{\normalsize cycles}
\end{region}

\begin{region}{338.9}
In the C program, modify the function main() as follows:
\end{region}

\begin{region}[pitch=10]{350.7}
\begin{asmlisting}[9]
 \aN{20} \aI{int} main()
 \aN{21} \{
 \aN{22}     \aC{// setup}
 \aN{23}     setPinModes(\aK{4});
 \aN{24}
 \aN{25}     \aI{unsigned} \aI{char} pinValues = \aK{0};
 \aN{26}     \aN{while} (\aK{1}) \{
 \aN{27}         \aC{// loop}
 \aN{28}         pinValues = pinValues \textasciicircum{} 4;
 \aN{29}         digitalWriteAll(pinValues);
 \aN{30}         delay\_cycles((\aK{500000000} / \aK{62.5} - \aK{7}) / \aK{6});
 \aN{31}     \}
 \aN{32} \}
\end{asmlisting}
\end{region}

\begin{region}{496.2}
\subheading{Task}
\end{region}

\begin{region}[pitch=13.3]{509.5}
\begin{numlist}
\item Make the LED blink at 1\,Hz.
\item Run \code{diff led_pin2_simple.hex ../asm/led_pin2_delay4.hex} and
  convince yourself that the machine code is identical to what we wrote in
  assembly language in Step~2.
\end{numlist}
\end{region}

\end{document}
